\begin{resumo}[Abstract]
 \begin{otherlanguage*}{english}
Machine learning models in production degrade over time due to data drift,
a change in the underlying data distribution that silently compromises the
validity of a trained model. While detecting this phenomenon is a
relatively mature problem, deciding how to respond to a detected drift
remains rarely investigated systematically, particularly regarding the
interaction between the drift's temporal pattern, its severity, and the
monitored classifier. This work proposes a controlled factorial experiment,
in a stream learning setting, to evaluate the effectiveness of four
retraining strategies (full reset, sliding window, incremental retraining,
and historical model reactivation), plus a no-retraining control condition,
under four temporal drift patterns, three severity levels, and three
classifiers with distinct mechanics (Naive Bayes, Random Forest, and
XGBoost). Severity is operationalized through a closed-form Hellinger
distance measure, computed both from the known generating distributions and
as estimated by the adopted detector (HDDDM), allowing this estimate to be
validated against the true value. As an exploratory extension, the work
also investigates whether observable signals from the detector allow
sketching a preliminary strategy-selection guideline applicable to MLOps
pipelines, where the true drift type is never directly known.
   \vspace{\onelineskip}

\noindent \textbf{Key-words}: data drift. concept drift. stream learning.
model retraining. drift detection.
 \end{otherlanguage*}
\end{resumo}